% ============================================================================
%  Chapter 18 -- The Processor                                   (Lecture 18)
% ============================================================================
\chapter{The Processor}
\label{ch:processor}

\epigraphline{One Ring to rule them all\ldots}
  {J.\,R.\,R.\ Tolkien, \emph{The Lord of the Rings}, 1954--1955}

\begin{chapgoals}
By the end of this chapter you will be able to
\begin{itemize}
\item name the components of a processor: control unit, ALU, registers, links to memory and to input/output;
\item distinguish internal registers from machine registers, and general registers from special ones;
\item explain how a stack works and how \riscv implements push and pop;
\item compare the accumulator, stack and general-register models on the same computation;
\item describe \riscv as a load-store, general-register processor, and read the name RV32IMAC;
\item map every element of the general model onto the processor of the \prg kit.
\end{itemize}
\end{chapgoals}
% Generated by tools/textbook/export_listings.js from tools/decks/mastery.js and tools/concepts/registry.json: do not edit.
\begin{chapstatus}{Learning status of this chapter}
\masterychip{now}{UNDERSTAND NOW}\enspace accumulator, stack, registers; the processor model; the registers of RISC-V; the calling convention; the stack; the RV32 processor.\par\smallskip
\masterychip{use}{USE NOW}\enspace the complete program, the PRG32 tools (used, not studied, in this course).\par\smallskip
\masterychip{later}{UNDERSTAND LATER}\enspace interrupts and traps (Chapter~\ref{ch:instructions}).
\end{chapstatus}
\begin{chapstatus}{Roadmap of this chapter}
By the end of the chapter you can:
\begin{enumerate}[nosep,leftmargin=2em]
\item follow an instruction through the blocks of an RV32 processor
\item use the stack and the calling convention correctly
\item compare accumulator, stack and register machines
\end{enumerate}
It builds on: only the setup of Chapter~\ref{ch:setup}. Next: Chapter~\ref{ch:instructions}.
\end{chapstatus}


The reader has programmed a processor since Chapter~\ref{ch:firstcontact}.
This chapter presents the general model of a processor, of which every real
one is a variant, and uses \riscv as the reference: every abstract element
has a name that the reader already knows.

% ----------------------------------------------------------------------------
\section{Warm-up on the machine}

\asmfile[caption={One line of a high-level language: \code{z = x + y}.},label=lst:cpu:sum]{code/l18/c03.S}

Seven lines for one sum (Listing~\ref{lst:cpu:sum}). Memory is touched only
by \ins{lw} and \ins{sw}, and the arithmetic happens only between registers:
it is the cycle ``load, modify, store'' of Chapter~\ref{ch:firstcontact}, now
seen as a design choice. Other processors do the same job with three
instructions, or with one. Why is \riscv made in this way? To answer, a
model of the processor is needed.

% ----------------------------------------------------------------------------
\section{The processor}

\subsection{The computer and its processor}

An electronic computer is a numeric, automatic, stored-program system built
with electronic circuits (Chapter~\ref{ch:boole1}), with three subsystems:
processor, main memory and input/output, connected by buses. The processor
can execute only a fixed set of instructions, from some tens to a few hundred,
more or less complex. The \term{machine language} is the set of the strings
of bits that encode the instructions which the circuits of a processor can
execute. Instructions request transfers of information, inside the computer
and to or from the I/O devices, and arithmetic and logic operations. The data
and the instructions of a program are stored in the main memory.

RV32I, the base of \riscv, has about forty instructions; with the M, A and C
extensions of the kit, about a hundred encodings. The Motorola 68000 had
fifty-six types of instruction, with many variants of each.

\subsection{A reference model}

Figure~\ref{fig:cpu:model} is the reference model of a processor. Its
elements are the following.

\begin{itemize}
\item The \term{control network}, or control unit, which produces the commands
  for all the rest.
\item The \term{program counter} PC, which holds the address of the next
  instruction, and the \term{decoding register} D, which holds the instruction
  being executed.
\item The \term{machine registers} $R_0,\dots,R_{n-1}$.
\item The \term{ALU}, with its input registers $O_1$ and $O_2$ and its output
  register $U$.
\item The link to memory: the address register MA and the buffer register MB
  of Chapter~\ref{ch:memory}.
\item The link to input/output: an input register RI and an output register RU.
\item An \term{internal bus} that joins them all.
\end{itemize}

\begin{figure}[htb]
\centering
\begin{tikzpicture}
  \node[draw=prgsteel,dashed,rounded corners,minimum width=116mm,minimum height=62mm,label={[lbl,anchor=north west]north west:CPU}] (cpu) at (5.3,-2.25) {};
  \node[blk,fill=prgsteel,text=white,minimum width=20mm,minimum height=17mm] (cn) at (1.0,-0.6) {control\\network};
  \node[blkred,minimum width=9mm,minimum height=7mm] (pc) at (0.5,-2.2) {PC};
  \node[blkred,minimum width=9mm,minimum height=7mm] (d) at (1.5,-2.2) {D};
  \foreach \i/\r in {0/{$R_0$},1/{$R_1$},2/{$\cdots$},3/{$R_{n-1}$}} \node[blkacc,minimum width=18mm,minimum height=5.5mm] (r\i) at (4.0,-\i*0.65+0.1) {\r};
  \node[blkvio,minimum width=11mm,minimum height=6mm] (o1) at (6.9,0.1) {$O_1$};
  \node[blkvio,minimum width=11mm,minimum height=6mm] (o2) at (8.3,0.1) {$O_2$};
  \node[blkvio,minimum width=25mm,minimum height=9mm] (alu) at (7.6,-0.85) {\textbf{ALU}};
  \node[blkvio,minimum width=11mm,minimum height=6mm] (u) at (7.6,-1.9) {$U$};
  \draw[prgamber,line width=3pt] (-0.2,-2.95) -- (10.8,-2.95) node[pos=0.93,above,lbl,text=prgamber]{internal bus};
  \node[blkgrn,minimum width=11mm,minimum height=7mm] (ma) at (1.0,-4.1) {MA};
  \node[blkgrn,minimum width=11mm,minimum height=7mm] (mb) at (2.6,-4.1) {MB};
  \node[blkamb,minimum width=11mm,minimum height=7mm] (ri) at (7.0,-4.1) {RI};
  \node[blkamb,minimum width=11mm,minimum height=7mm] (ru) at (8.6,-4.1) {RU};
  \foreach \n in {pc,d,u} \draw[<->,thick] (\n.south) -- (\n.south|-0,-2.95);
  \draw[<->,thick] (r3.south) -- (r3.south|-0,-2.95);
  \foreach \n in {ma,mb,ri,ru} \draw[<->,thick] (\n.north) -- (\n.north|-0,-2.95);
  \draw[->,thick] (o1) -- (o1|-alu.north); \draw[->,thick] (o2) -- (o2|-alu.north); \draw[->,thick] (alu) -- (u);
  \draw[<->,bus] (1.8,-4.45) -- (1.8,-6.0) node[below,lbl]{to the memory};
  \draw[<->,bus] (7.8,-4.45) -- (7.8,-6.0) node[below,lbl]{to input/output};
\end{tikzpicture}
\caption{The reference model of a processor.}
\label{fig:cpu:model}
\end{figure}

Although they are based on the same logical scheme, and are therefore made of
rather similar components, the processors on the market can be very different
from one another. Table~\ref{tab:cpu:map} puts the model and \riscv side by
side. The two rows marked ``absent'' are deliberate simplifications of
\riscv, discussed below and in Chapter~\ref{ch:memory}. The MC68000 of the
previous editions of this course had all the elements of the model: eight
data registers and eight address registers, a status register with condition
codes, and separate circuitry for interrupts.

\begin{table}[htb]
\centering\small
\caption{The model and \riscv.}\label{tab:cpu:map}
\begin{tabularx}{\textwidth}{@{}P{0.3\textwidth}LL@{}}
\toprule
\textbf{Element of the model} & \textbf{In \riscv} & \textbf{Seen by the programmer?}\\
\midrule
PC, program counter & \reg{pc}, 32 bits & indirectly: branches, jumps and \ins{auipc} use it\\
D, decoding register & the instruction register & no\\
$R_0\ldots R_{n-1}$, machine registers & \reg{x0}\ldots\reg{x31}, thirty-two registers of 32 bits & yes: \reg{zero}, \reg{ra}, \reg{sp}, \reg{t0}, \reg{a0}, \reg{s0}, \ldots\\
$O_1$, $O_2$, $U$, registers of the ALU & the inputs and the output of the ALU & no\\
MA, MB, link to memory & the address and the datum of a load or store & no\\
RI, RU, link to input/output & absent: devices are memory-mapped & --\\
status register with flags & absent: no condition codes & --\\
special registers & the control and status registers (CSR) & only by privileged software\\
\bottomrule
\end{tabularx}
\end{table}

\subsection{Internal registers and machine registers}

\begin{definition}{Internal and machine registers}
The \term{internal registers} are not directly visible to the programmer, and
are necessary for the processor to work: PC, MA, MB, D, $O_1$, $O_2$, $U$.
The \term{machine registers} define the programming model; they are divided
into \term{general} registers and \term{special} registers.
\end{definition}

A processor handles three kinds of information: instructions, memory
addresses and data. The size of MA and of the PC depends on the size of the
addresses; the size of the registers of the ALU depends on the size of the
data. In RV32 all three are 32 bits wide: one size for registers, addresses
and instructions. The ``32'' of the name is the width of the registers: RV64
has registers and addresses of 64 bits, and the same instructions of 32 bits.

\subsection{How real processors differ from the model}

In general, real processors may have a different and usually richer set of
general registers; their machine registers may be specific memory locations;
the I/O section may be missing altogether and be implemented by dedicated
chips; and very simple arithmetic units can be extended by dedicated chips
that improve the ALU, as the floating-point coprocessors 68881 for the 68000
family and 8087 for the 8086 did.

\riscv on the \prg kit has thirty-two general registers in a bank inside the
processor; no I/O section, since the peripheral controllers are on the same
chip and are reached as memory; optional \emph{extensions} instead of extra
chips, of which M adds multiplication and division; and no F extension, so
that floating point would be done in software (Chapter~\ref{ch:reals}).

\subsection{Words or bytes: the length of machine words}

In a \emph{word processor} all the machine words have the same length, sized
on the longest piece of information; this is awkward with data of
intrinsically variable nature. Every \riscv instruction is 32 bits long (16
in the compressed form). In a \emph{character}, or byte, processor the length
of machine words is variable, a multiple of a byte, which suits the treatment
of strings: an Intel x86 instruction takes from 1 to 15 bytes, and an MC68000
instruction from one to five words of 16 bits. A characteristic parameter of
a processor is the number of bits of its registers, the parallelism of its
internal transfers: 8, 16, 32 or 64.

\subsection{User state and supervisor state}

A processor is in \term{supervisor state} when it executes the operating
system, and may then execute every instruction of its set; it is in \term{user
state} when it executes a user program, and may then execute everything
except the \term{privileged instructions}.

\riscv defines three \term{privilege modes}. Machine mode (M), the most
privileged and the only mandatory one, runs the firmware. Supervisor mode
(S) runs the kernel of an operating system such as Linux. User mode (U) runs
applications. The ESP32-C6 implements M and U. The passage from one mode to
another happens through traps, the subject of Chapter~\ref{ch:instructions}.

% ----------------------------------------------------------------------------
\section{Registers}

\subsection{Kinds of registers}

A register can be specialised, or can do several jobs. Specialised registers
need specialised operation codes for the same operation on different banks;
concentrating several functions in the same registers gives more internal
complexity and more flexibility of use. The machine registers, different in
number and often grouped in banks, can be organised as an addressable memory
(a bank $R$, of which $R_i$ is a register) or as autonomous registers, each
with its own name. Some processors placed their general registers in the
first locations of memory: the ``page zero'' of the 6502.

\riscv has one bank of thirty-two registers, addressed by 5 bits inside the
instruction; the names \reg{t0}, \reg{a0}, \reg{sp} are only conventions.

\subsection{The functions of registers}

\begin{description}[style=nextline,leftmargin=1.5em]
\item[Transit of data.] From memory to a register, from register to register,
  from a register to memory: load, move, store.
\item[Accumulation of results.] Needed for arithmetic and logic:
  $R:=\mathrm{op}\,R$, or $R:=R\ \mathrm{op}\ X$.
\item[Addressing.] Finding the address of an operand: index registers, base
  registers.
\item[Stack pointers.] Managing a stack implemented in memory.
\item[Indicators, or flags.] Binary information about events: zero, sign,
  overflow, carry, illegal addresses or codes. Produced during the execution
  of an instruction, they can condition the following ones.
\end{description}

In \riscv the first four functions are all served by the general registers;
the fifth does not exist as a register.

\subsection{The thirty-two registers of RISC-V}

\begin{table}[htb]
\centering\small
\caption{The registers of \riscv and their conventional roles.}\label{tab:cpu:regs}
\begin{tabular}{@{}lll@{\qquad}lll@{}}
\toprule
\textbf{Register} & \textbf{Name} & \textbf{Function} & \textbf{Register} & \textbf{Name} & \textbf{Function}\\
\midrule
\reg{x0} & \reg{zero} & always 0 & \reg{x8}--\reg{x9} & \reg{s0}--\reg{s1} & saved\\
\reg{x1} & \reg{ra} & return address & \reg{x10}--\reg{x11} & \reg{a0}--\reg{a1} & arguments, results\\
\reg{x2} & \reg{sp} & stack pointer & \reg{x12}--\reg{x17} & \reg{a2}--\reg{a7} & arguments\\
\reg{x3} & \reg{gp} & global pointer & \reg{x18}--\reg{x27} & \reg{s2}--\reg{s11} & saved\\
\reg{x4} & \reg{tp} & thread pointer & \reg{x28}--\reg{x31} & \reg{t3}--\reg{t6} & temporaries\\
\reg{x5}--\reg{x7} & \reg{t0}--\reg{t2} & temporaries & \reg{pc} & -- & program counter\\
\bottomrule
\end{tabular}
\end{table}

The hardware treats \reg{x1} to \reg{x31} identically, and only \reg{x0} is
special (Table~\ref{tab:cpu:regs}); the program counter is not in the bank.
The roles come from the calling convention, met in Chapter~\ref{ch:firstcontact},
whose rule is this. A call may destroy \reg{ra} and the \reg{t} and \reg{a}
registers: they are \emph{caller-saved}, and a caller that still needs one saves
it before the call. A call must return the \reg{s} registers and \reg{sp}
unchanged: they are \emph{callee-saved}, and a function that uses one saves it
in its stack frame and restores it before \ins{ret}. Appendix~\ref{app:riscv}
gives the complete convention. The MC68000, instead, separated
the data registers D0--D7 from the address registers A0--A7, with A7 as stack
pointer.

\subsection{Flags: the status register, and its absence}

Many processors have a \term{status register}: bits, called flags, accessible
from the program. The usual ones are Z (zero), N (negative), V (overflow) and
C (carry), sometimes with a parity bit. An instruction sets them, and a later
conditional branch tests them. The MC68000 has its condition code register,
and x86 its flags register.

\riscv has no flags. A branch compares two registers and jumps, in one
instruction: \ins{beq}, \ins{bne}, \ins{blt}, \ins{bge}, \ins{bltu},
\ins{bgeu}. The instructions \ins{slt} and \ins{sltu} write the result of a
comparison in a general register. Carry and overflow are computed when they
are needed (Chapter~\ref{ch:integers}). There is less hidden state, which
makes pipelines and interrupts simpler.

\subsection{Special registers of RISC-V: the CSRs}

The special registers of \riscv are the \term{control and status registers}
(CSR): a separate bank of up to 4\,096 registers, read and written only with
the instructions \ins{csrrw}, \ins{csrrs} and \ins{csrrc}, and only by
privileged code (Table~\ref{tab:cpu:csr}). They are the ``registers of the
interrupts'' of the general model, and the subject of
Chapter~\ref{ch:instructions}.

\begin{table}[htb]
\centering\small
\caption{Some machine-mode CSRs.}\label{tab:cpu:csr}
\begin{tabularx}{\textwidth}{@{}llL@{}}
\toprule
\textbf{Register} & \textbf{Name} & \textbf{What it holds}\\
\midrule
\reg{mstatus} & machine status & the global interrupt enable and the previous privilege mode\\
\reg{mtvec} & trap vector & the address of the routine that handles interrupts and exceptions\\
\reg{mepc} & exception program counter & the address of the instruction that was interrupted\\
\reg{mcause} & trap cause & a code that says why the trap happened\\
\reg{mie}, \reg{mip} & interrupt enable, interrupt pending & one bit per source of interrupt\\
\bottomrule
\end{tabularx}
\end{table}

% ----------------------------------------------------------------------------
\section{The stack}

\subsection{Last in, first out}

\begin{definition}{Stack}
A \term{stack} is an abstract data structure with two operations: \emph{push}
stores an element as the last of the pile, and \emph{pop} takes the last
element away. The last value pushed is the first one popped (LIFO).
\end{definition}

A stack works always and only on its top: there is no need to say where to
store or where to fetch. It is used to set aside the contents of registers
that must be freed to carry out a task, and to restore them afterwards: this
happens in subprograms and in interrupts. It is implemented with a pointer,
in memory or in a dedicated register, which indicates at every instant the
position of the last element. Operations on a stack should check for overflow
on a push and for an empty stack on a pop.

\subsection{Push and pop in RISC-V}

In \riscv the stack lives in memory, grows towards the low addresses, and the
register \reg{sp} points to its top (Figure~\ref{fig:cpu:stack}). There are
no instructions called push and pop: each is two ordinary instructions.

\begin{figure}[htb]
\centering
\begin{minipage}{0.42\textwidth}\centering
\begin{tikzpicture}[x=0.7cm,y=0.62cm,font=\ttfamily\small]
  \foreach \i/\v/\f in {0/(free)/prgpanel, 1/(free)/prgpanel, 2/-28/prgamber!60, 3/-17/prgaccent!18, 4/739/prgaccent!18, 5/43/prgaccent!18}{
    \draw[prgsteel,fill=\f] (0,-\i) rectangle (2.6,-\i-1);
    \node at (1.3,-\i-0.5) {\v};
  }
  \draw[->,very thick,prgred] (-1.1,-2.5) node[left,font=\ttfamily\bfseries]{sp} -- (-0.05,-2.5);
  \node[lbl,anchor=west,text=prgcomment] at (2.7,-0.5) {low addresses};
  \node[lbl,anchor=west,text=prgcomment] at (2.7,-5.5) {bottom};
\end{tikzpicture}
\end{minipage}
\begin{minipage}{0.5\textwidth}
\begin{asm}[numbers=none,xleftmargin=4pt,basicstyle=\ttfamily\footnotesize]
# push t0
addi sp, sp, -4      # make room
sw   t0, 0(sp)       # store on the new top

# pop t0
lw   t0, 0(sp)       # read the top
addi sp, sp, 4       # release it
\end{asm}
\end{minipage}
\caption{A stack in memory, and the two operations in \riscv.}
\label{fig:cpu:stack}
\end{figure}

In practice a \riscv program does not push and pop single values: a
\emph{stack frame} is allocated once per function, sixteen bytes at a time,
and is filled with ordinary stores, as every function of \vermone that makes
a call has done. Other processors have dedicated forms: on the 68000,
\code{MOVE (SP)+, D0} pops with a post-increment.

% ----------------------------------------------------------------------------
\section{Processor models}

Processors are classified according to where the operands of an operation
live.

\subsection{Four architectures}

\begin{description}[style=nextline,leftmargin=1.5em]
\item[General registers.] The processor has a set of registers that can be
  used interchangeably.
\item[Accumulator.] All the arithmetic, logic and comparison instructions have
  an implicit operand: one special register, the accumulator.
\item[Stack.] There are neither accumulators nor general registers, only the
  stack pointer; the instructions work on the data on top of the stack.
\item[Orthogonal.] A single bank of general registers, some of which serve,
  when needed, as stack pointer or as other special registers.
\end{description}

\subsection{\texttt{z = x + y} in four processors}

\begin{table}[htb]
\centering\small
\caption{The same computation in four models.}\label{tab:cpu:models}
\begin{tabularx}{\textwidth}{@{}P{0.21\textwidth}P{0.2\textwidth}Lc@{}}
\toprule
\textbf{Model} & \textbf{Example} & \textbf{Program} & \textbf{Accesses}\\
\midrule
accumulator & Motorola 6809, MOS 6502 & \code{LDA x; ADD y; STA z} & 3\\
stack & Java virtual machine, WebAssembly & \code{PUSH x; PUSH y; ADD; POP z} & 3 + stack\\
general registers, memory operands & MC68000, Intel x86 & \code{MOVE x,D0; ADD y,D0; MOVE D0,z} & 3\\
general registers, load-store & \riscv & \code{lw t1,x; lw t2,y; add t3,t1,t2; sw t3,z} & 3\\
\bottomrule
\end{tabularx}
\end{table}

The accesses to memory are the same three everywhere
(Table~\ref{tab:cpu:models}): what changes is how many instructions there are,
and how complex each one is. In the register-memory model an instruction
such as \code{ADD y, D0} reads memory \emph{and} computes; in the load-store
model every instruction does one thing. Stack machines with the stack in
memory are generally slower than general-register machines. (In the table,
\code{lw t1, x} is a shorthand that the assembler accepts for the pair
\ins{la} and \ins{lw} of the warm-up.)

\subsection{RISC-V: a load-store, general-register processor}

\begin{prgconcept}[Load, compute, store]
Only loads and stores touch memory; every other instruction works from
register to register. It is the rhythm of every \riscv program, from the
first counter to the last game.
\end{prgconcept}

The reasons are four. Instructions that work on registers are faster than
those that work on memory locations, so that the most used data are kept in
registers. An instruction has three operands, two sources and a destination,
so that no value is destroyed unless the programmer wants it to be. Each
instruction does one simple thing and has a fixed size: it is easy to decode,
and easy to execute in a pipeline. This is the meaning of RISC, \emph{reduced
instruction set computer}: the forms that go from memory to a register and
from a register to memory, which exist in CISC processors, are exactly what
RISC gives up.

% ----------------------------------------------------------------------------
\section{The ALU and the whole picture}

\subsection{The arithmetic logic unit}

The ALU of Chapter~\ref{ch:combinational}, seen from the processor,
supervises the arithmetic and logic operations between operands. In the model
it has two input registers, $O_1$ and $O_2$, and an output register $U$. It
is physically included in the processor in elementary architectures and is an
independent apparatus in more complex ones; a widespread solution is a reduced
arithmetic unit inside the processor with independent apparatus for the more
sophisticated operations. In \riscv that apparatus is the extensions: M for
multiplication and division, F and D for floating point. The operations
available depend on the types of data and on their representation: logical
values, strings of characters, pointers, and numbers in the encodings of
Chapters \ref{ch:integers} and~\ref{ch:reals}.

\subsection{The architecture of an RV32 processor}

Figure~\ref{fig:cpu:datapath} reduces to its blocks the classic single-cycle
\riscv datapath of \citet{patterson2020}. The PC gives the address of the
instruction; the instruction is fetched from memory; the decoder separates its
fields and the control unit produces the commands; the register file supplies
the operands; the ALU computes a result or forms an address; the data memory
is read or written; and the result goes back to the register file. The
instruction memory and the data memory are the same address space, drawn
twice for clarity.

\begin{figure}[htb]
\centering
\begin{tikzpicture}[node distance=4mm]
  \node[blkred,text width=12mm,minimum height=16mm] (pc) {\textbf{PC}};
  \node[blkvio,text width=19mm,minimum height=16mm,right=of pc] (im) {\textbf{instruction memory}\\\scriptsize fetch};
  \node[blk,fill=prgsteel,text=white,text width=20mm,minimum height=16mm,right=of im] (dec) {\textbf{decoder and control}\\\scriptsize fields, commands};
  \node[blkacc,text width=19mm,minimum height=16mm,right=of dec] (rf) {\textbf{register file}\\\scriptsize \texttt{x0}--\texttt{x31}};
  \node[blkamb,text width=16mm,minimum height=16mm,right=of rf] (alu) {\textbf{ALU}\\\scriptsize compute, or form an address};
  \node[blkvio,text width=16mm,minimum height=16mm,right=of alu] (dm) {\textbf{data memory}\\\scriptsize load, store};
  \foreach \a/\b in {pc/im,im/dec,dec/rf,rf/alu,alu/dm} \draw[->,bus] (\a) -- (\b);
  \draw[->,bus,prgred,rounded corners=2mm] (dm.south) -- ++(0,-8mm) -| (rf.south) node[pos=0.25,below,lbl,text=prgred]{write back the result};
  \draw[->,bus,prgsteel,rounded corners=2mm] (alu.north) -- ++(0,8mm) -| (pc.north) node[pos=0.25,above,lbl]{next PC: $PC+4$ or a branch target};
\end{tikzpicture}
\caption{The architecture of an RV32 processor, reduced to its blocks.}
\label{fig:cpu:datapath}
\end{figure}

\subsection{The processor of the PRG32 kit}

\begin{table}[htb]
\centering\small
\caption{The ESP32-C6 (from the datasheet of Espressif).}\label{tab:cpu:c6}
\begin{tabularx}{\textwidth}{@{}lL@{}}
\toprule
\textbf{Feature} & \textbf{ESP32-C6}\\
\midrule
Instruction set & RV32IMAC: base integer set, multiply and divide, atomics, compressed instructions\\
Main core & one 32-bit \riscv core, up to 160 MHz, four-stage pipeline\\
Second core & a low-power \riscv core, up to 20 MHz, for tasks that run while the main core sleeps\\
Privilege modes & machine and user\\
Floating point & none: no F or D extension\\
Memory & 512 KB of SRAM, 320 KB of ROM, external flash\\
On the same chip & Wi-Fi 6, Bluetooth, Thread and Zigbee radios; SPI, I2S, GPIO and timer peripherals\\
\bottomrule
\end{tabularx}
\end{table}

The chip of Table~\ref{tab:cpu:c6} is a complete computer on a few square
millimetres that costs about two euros: processor, memory and input/output
are all there, as in the model with which the chapter began.

% ----------------------------------------------------------------------------
% Generated by tools/textbook/export_listings.js from decks/plan.js and code/l18.S: do not edit.
\section{The program of this chapter}\label{sec:lab:l18}

The listings of this chapter are parts of one program, \fname{code/l18.S} in the
repository of the course, printed in full in Listing~\ref{lst:lab:l18}. The
automatic tests of the course material build it and run it on the headless
runner of \prgqt{}, with scripted input, and check the frame;
Figure~\ref{fig:lab:l18} is the frame that the test wrote.

\begin{description}[style=nextline,leftmargin=1.5em]
\item[On the screen.] Three bars, eight pixels per unit: x = 5 (white), y = 7 (white) and z = 12 (green), computed by the seven instructions of the slide.
\item[What to press.] Nothing: the program runs by itself.
\end{description}

\begin{figure}[htb]
\centering
\includegraphics[width=0.5\textwidth]{lab_l18.png}
\caption{\fname{l18.S} on \prgqt{}.}\label{fig:lab:l18}
\end{figure}

\asmfile[caption={The complete program of the chapter, \fname{code/l18.S}.},label=lst:lab:l18]{code/l18/program.S}

To run it on \prgqt{} the commands are those of Listing~\ref{lst:lab:l18:run};
Section~\ref{sec:prg:testing} gives the same steps for QEMU and for the board.

\begin{shell}[caption={Building \fname{l18.S} and running it on \prgqt{}.},label=lst:lab:l18:run]
# 1. start PRG32-QT on this computer (it listens on port 8080)

# 2. in the PRG32 folder, with ESP-IDF active: build the cartridge
python3 -m prg32 cartridge build ~/arc12/code/l18.S --portable \
    --entry-prefix l18 --name l18 --out l18.prg32

# 3. send it to PRG32-QT and start it
python3 -m prg32 esp32c6 upload-and-run l18.prg32 \
    --url http://localhost:8080 --slot cart0

# 4. look, press, compare with the previous slide
\end{shell}


% ----------------------------------------------------------------------------
% Generated by tools/textbook/export_listings.js from tools/concepts/assessments.json: do not edit.
\begin{chapstatus}{The thread of this chapter}
The ra saved in Lecture 02 is now explained: the stack and the calling convention are rules that let functions written by different people call each other. Next: how an instruction is encoded and executed, and how the processor is interrupted.
\par\medskip\textbf{Exit ticket} (two to four minutes, alone, on paper):
\begin{enumerate}[nosep,leftmargin=2em]
\item Follow add t2, t0, t1 through the block diagram of an RV32 processor: name, in order, the blocks it uses and what each does. Which block does it not use?
\item A function begins with addi sp, sp, -16 and sw ra, 12(sp). Where is ra now, and what must the function do before ret?
\end{enumerate}
\end{chapstatus}


\section{Check your understanding}

\begin{quiz}{Which register of the general model has no counterpart in \riscv?}
\item the program counter
\item the status register with the flags
\item the stack pointer
\item the instruction register
\end{quiz}

\begin{quiz}{In \riscv, which instructions can read or write the main memory?}
\item all of them
\item \ins{add} and \ins{sub}
\item only loads and stores
\item only jumps
\end{quiz}

\begin{quiz}{After \code{addi sp, sp, -4} and \code{sw t0, 0(sp)}, what has happened?}
\item \reg{t0} was popped from the stack
\item \reg{t0} was pushed on the stack
\item the stack was emptied
\item \reg{sp} was saved in \reg{t0}
\end{quiz}

\begin{quiz}{In an accumulator machine, the instruction \code{ADD y} means:}
\item $y := y + \text{accumulator}$
\item $\text{accumulator} := \text{accumulator} + [y]$
\item push $y$, then add
\item $y := \text{accumulator}$
\end{quiz}

\section{Exercises}

\begin{exercise}[paper first, then \prg]
\begin{enumerate}[nosep]
\item Write $z=(a+b)-(c+d)$ for an accumulator machine, for a stack machine
  and for \riscv; count the instructions and the memory accesses.
\item Write the push and the pop of two registers, \reg{t0} and \reg{t1},
  keeping \reg{sp} a multiple of 16 as the calling convention requires.
\item For each register used in your last \prg program, say which function it
  serves: transit, accumulation, addressing, stack.
\item In \prgqt, single-step a call and its return: watch \reg{pc}, \reg{ra}
  and \reg{sp}, and the word saved on the stack.
\item Find in the datasheet of the processor of your telephone or laptop:
  instruction set, number of cores, clock, privilege levels.
\end{enumerate}
\expected{Accumulator: \code{LDA a; ADD b; STA tmp; LDA c; ADD d; STA tmp2;
LDA tmp; SUB tmp2; STA z}. Stack: \code{PUSH a; PUSH b; ADD; PUSH c; PUSH d;
ADD; SUB; POP z}. \riscv: four loads, two \ins{add}, one \ins{sub}, one
store. In \prgqt: \reg{sp} drops by 16 at the prologue and returns at the
epilogue; \reg{ra} holds the address after the call.}
\end{exercise}

\section{Glossary}

\begin{glossarytable}
control unit & unit\`a di controllo\\
program counter & contatore di programma\\
machine register & registro di macchina\\
general / special register & registro generico / speciale\\
instruction set & insieme (set) di istruzioni\\
accumulator & accumulatore\\
stack pointer & puntatore allo stack\\
status register, flags & registro di stato, indicatori\\
supervisor / user state & stato supervisore / utente\\
privileged instruction & istruzione privilegiata\\
\end{glossarytable}

\section{Conclusions and reading}

The processor commands and controls, relying on the memory and on the
input/output systems. Registers are storage locations that may be physically
inside or outside the processor, and the kinds of registers and the way in
which they are used determine the model of the processor: accumulator, stack,
general registers. The ALU carries out the operations between operands.
\riscv is a load-store, general-register processor with no flags and no I/O
instructions: few elements, each one simple.

\paragraph{Reading.} \citet{patterson2020}, chapters 2 and~4; the \riscv
instruction set manual, volume~I \citep{riscvisa}; \citet{stallings2021}, the
chapters on the structure and function of the processor.
Appendix~\ref{app:riscv} is the complete reference.
